\subsection{Secondary terminal intervals with their entrances}
\label{rw:anchors:secondary}

Let $R$ be the regular, wedge and regular-secondary prefix, $T$ its
finite-saddle or weak terminal interval, and $O$ their actual open
overlap. We use the common $\mathcal K$ of
\eqref{rw:secondary:complete-log}, the fixed-fibre forcing of
Section~11 and the seam convention of Section~\ref{rw:anchors:primary}.

By Proposition~\ref{rw:regular-chain}, the chain has at most one
simple zero, after which
\begin{equation}\label{rw:secondary:anchor-entry}
 \bar aC<0,\qquad
 \operatorname{sign}\mathcal K=-\operatorname{sign}C,\qquad
 \operatorname{sign}\partial_{\rm out}\mathcal K
                                      =-\operatorname{sign}C .
\end{equation}
Here $C=\bar d/\lambda$ with $\lambda>0$ fixed on the fibre.
Lemma~\ref{rw:secondary:chain-monotone} gives the same departure
from a supplied signed entrance; the positive column quotient on a
compact regular-secondary overlap gives $|\bar a|\le K_1|C|$.

\begin{proposition}[The finite secondary saddle with an entrance]
\label{rw:secondary:saddle-anchor}
Let $T$ be the finite-saddle interval of
Proposition~\ref{rw:secondary:saddle}, with its fixed-fibre relative
forcing. Without a prescribed entrance it has
$Z(\mathcal K;T)\le2$.
If the preceding chain has a zero, then its remaining terminal part
has at most one zero of $\mathcal K$. The same improvement holds
without a preceding zero when $O$ satisfies
\begin{equation}\label{rw:secondary:saddle-signed-entry}
 \bar aC<0,\qquad
 \operatorname{sign}\mathcal K|_O=-\operatorname{sign}C .
\end{equation}
Thus $Z(\mathcal K;R\cup T)\le2$ without a specified entrance,
and the terminal union after an entrance satisfying
\eqref{rw:secondary:saddle-signed-entry} has at most one zero
if its preceding chain has no zero.
All counts include multiplicities and use the actual passing and
positive-central restrictions.
\end{proposition}

\begin{proof}
First separate the lower-double overlay. On a compact regular overlap
before the saddle, the upper logarithms and the relative lower term
are bounded, whereas \eqref{rw:secondary:lower-constant} has arbitrarily
large modulus and sign $-\operatorname{sign}\bar a$ near
$\bar a=\pm2\sqrt{\bar B}$. Under
\eqref{rw:secondary:saddle-signed-entry} this is
$\operatorname{sign}C$, opposite to the entrance sign, so the anchored
case is excluded. The value-domination argument of
Proposition~\ref{rw:secondary:saddle} already gives at most one
zero in the overlay, all crossings having the same nonzero actual-ratio
sign. This count survives restriction by every lower whole-path test.

On the remaining compact regular-lower block the signed comparison
maps extend along the saddle quadrant to $t=0$, with
\begin{equation}\label{rw:secondary:anchor-saddle-log}
 \mathcal K=K_c+\epsilon\kappa\log t+\mathcal E(t),
 \qquad \mathcal E(t)\longrightarrow0 .
\end{equation}
Here $\kappa=r-1$ is the actual ratio difference. The two central outputs
need not both be positive on this whole continuation.

To determine the endpoint sign of this continuation, we need the sign
of the actual ratio difference $\kappa$.
The overlap cone and the transported forcing give
\[
 |\bar a|\le K_1|C|,\qquad
 |\nu|\le\eta_\epsilon(|\bar a|+|C|)
                    \le (K_1+1)\eta_\epsilon|C|.
\]
Near $C=0$, identity \eqref{rw:secondary:saddle-ratio} now reads
\[
 \kappa=-\alpha_s C+\mathcal R_s,\qquad
 |\mathcal R_s|\le
 C_0\bigl(|C|+\epsilon\lambda K_1+
           \epsilon\lambda(K_1+1)\eta_\epsilon\bigr)|C|,
 \qquad \alpha_s>c_0>0 .
\]
Choose the $C$ neighborhood and then the scale restrictions so that
the coefficient of $|C|$ in this error is less than $c_0/2$.
Thus $\operatorname{sign}\kappa=-\operatorname{sign}C$ and
$|\kappa|\ge c_0|C|/2$. Away from this neighborhood the same sign
follows with a strict margin from
\[
 r_0-1=-\frac{2C}{\sqrt{C^2+4M}+C};
\]
the exact divided field is a uniformly small perturbation on the
compact positive-$M$ block. The case $C=0$ is excluded by
$\bar aC<0$. The retained coefficient error $\epsilon\lambda|\nu|$
is therefore small relative to $|C|$, which gives the ratio sign
uniformly even as $C$ tends to zero.

For $|\bar a|>K_0|\kappa|$ the same value-domination argument
already gives at most one zero. In the complementary cone the
unanchored estimate $Z(\mathcal K)\le2$ holds on the entire
saddle comparison continuation. Equation
\eqref{rw:secondary:anchor-saddle-log} gives its endpoint sign
$\operatorname{sign}C$, opposite to the sign at the fixed seam.
The opposite endpoint signs and Lemma~\ref{rw:anchors:calculus}
give one zero on the continuation and at most one on its physical
subinterval. The departure
\eqref{rw:secondary:anchor-entry} keeps the overlap zero-free after a
chain zero; adding that possible zero gives the union count.
\end{proof}

\begin{proposition}[The weak secondary end with an entrance]
\label{rw:secondary:weak-anchor}
Consider the weak row of Proposition~\ref{rw:secondary:real},
first with $C>0$. Without an entrance condition its terminal
interval has $Z(\mathcal K;T)\le2$.
If the preceding chain has a zero, its remaining terminal part
has at most one zero. This improvement also holds when there is
no chain zero but the open overlap satisfies
\begin{equation}\label{rw:secondary:weak-signed-entry}
 \mathcal K|_O<0,\qquad \partial_{\rm out}\mathcal K|_O<0 .
\end{equation}
Consequently the whole chain and terminal union has at most two
zeros of $K$; with the supplied entrance and no chain zero,
it has at most one. Reflection supplies the $C<0$ statement.
These conclusions permit either lower double root to shorten the
passing interval and permit positive-central truncation.
\end{proposition}

\begin{proof}
On the physical fibre Section~11 supplies
\[
 |\nu|\le C\delta_\epsilon(|\bar a|+|\bar d|),\qquad
 \bar d=\lambda C,
\]
from the outer closing zero. After a chain zero, or under the entrance,
$\mathcal K<0$ and $\partial_{\rm out}\mathcal K<0$ there.
The actual outgoing tag is
$P=P_s-\psi(t(u))$ with $\psi'>c$ and $t'>0$,
so outward motion decreases $u$ and
$\operatorname{sign}\partial_{\rm out}\mathcal K
=-\operatorname{sign}\mathcal M$.
Lemma~\ref{rw:anchors:calculus}, with the increasing numerator of
Proposition~\ref{rw:secondary:real}, leaves at most one simple tail
zero on the lower-passing, positive-central interval. Adding the
possible chain zero gives the assertions; reflection supplies $C<0$.
\end{proof}
